\section{Entropy transfer along the observation path}\label{cap:entropy}

The curvature estimates of Theorem~\ref{cap:true-curvature} control the
information revealed by the edge observations. We now use that control to
retain enough relative entropy until the posterior becomes a product measure.
The product inequality then gives a lower bound on the original heat-bath
Dirichlet form. All constants in this section may depend on fixed scale
cutoffs and fixed tolerances, but not on the test function or on $n$.

For a probability measure $\eta$ on the cube, write
\[
 \operatorname{Ent}_{\eta}(u)
 =\eta[u\log u]-\eta[u]\log\eta[u],\qquad
 \mathcal D_{\eta}(f)=\sum_{i=1}^n
       \eta\big[\operatorname{Var}_{\eta}(f\mid X_{-i})\big].
\]
Here $0\log0=0$, $\log_+u=\max\{\log u,0\}$, and $\mathcal D_0=\mathcal
D_\mu$. The interaction matrix is the GOE matrix $J$, with off-diagonal
variance $1/n$; its diagonal does not change $\mu(x)\propto\exp(x^{\mathsf
T}Jx/2)$.

\subsection{A reduction to bounded tests}\label{cap:capacity-reduction}

It is enough to prove the following estimate for every fixed $r>0$:
\begin{equation}\label{cap:restricted-energy}
 \mathcal D_0(f)\ge n^{-2/3-r}N L,
 \qquad N=\mu[f^2]=e^{-L},\qquad
 \nu=\frac{f^2\mu}{N},
\end{equation}
whenever $0\le f\le1$, $0<N$, $\mu(f>0)\le1/2$, and
\begin{equation}\label{cap:restricted-entropy}
 D(\nu\Vert\mu)\ge L/2,\qquad L\le Cn.
\end{equation}
The probability event on which this holds must be independent of $f$. The
observation argument is most effective for a bounded function whose
squared-density tilt has entropy comparable to its logarithmic inverse mass.
We first show that proving the energy estimate only for those tests is enough.
Equal height layers turn it into a capacity estimate for sets; dyadic layers
then control arbitrary half-supported functions. A median decomposition and an
entropy shift recover every real-valued test.
For this class, $\log2\le L$ and
$L/2\le D(\nu\Vert\mu)\le L$, while $d\nu/d\mu=f^2/N\le e^L$.

\begin{lemma}[Layers and capacity]\label{cap:layer-lemma}
Let $\eta$ be a strictly positive probability measure on a finite state
space, let $\mathcal E$ be a reversible Markov Dirichlet form, and suppose
$\min_x\eta(x)\ge e^{-M}$, with $M\ge\log2$. Suppose that, for some $A>0$,
\[
 \mathcal E(g)\ge A\eta[g^2]\log\frac1{\eta[g^2]}
\]
for every $0\le g\le1$ supported on a set of measure at most $1/2$ and
satisfying
\[
 D\!\left(\frac{g^2\eta}{\eta[g^2]}\middle\Vert\eta\right)
 \ge\frac12\log\frac1{\eta[g^2]},\qquad
 0<\log\frac1{\eta[g^2]}\le M.
\]
Then, with an absolute constant $C$,
\begin{equation}\label{cap:layer-lsi}
 \operatorname{Ent}_{\eta}(f^2)
 \le \frac{C(1+\log_+M)^2}{A}\,\mathcal E(f)
 \qquad\text{for every real-valued }f.
\end{equation}
\end{lemma}

\begin{proof}
Fix a nonempty set $E$ of measure $a\le1/2$ and a test $h$ equal to one
on $E$ and zero outside a set of measure at most $1/2$.  Clipping $h$ to
$[0,1]$ decreases its form.  Set
\[
 K=\left\lceil\log_2\!\left(1+
                  \frac{\log(1/a)}{\log2}\right)\right\rceil,
 \qquad
 g_i=K\big((h-i/K)_+\wedge K^{-1}\big),\quad 0\le i<K.
\]
Let $S_i=\eta(h>i/K)$ and $T_i=\eta(h\ge i/K)$. Some $i$ satisfies $T_{i+1}\ge
S_i^2$. Otherwise $T_1<S_0^2\le1/4$ and $T_{i+1}<T_i^2$, which would give
$a\le T_K<2^{-2^K}<a$. For this layer put $N_i=\eta[g_i^2]$. Then $a\le
T_{i+1}\le N_i\le S_i\le1/2$ and
\[
 D\!\left(\frac{g_i^2\eta}{N_i}\middle\Vert\eta\right)
 \ge\log(1/S_i)\ge\tfrac12\log(1/N_i).
\]
The first inequality follows by conditioning on its support. Also
$\log(1/N_i)\le\log(1/a)\le M$, and $\mathcal E(g_i)\le K^2\mathcal E(h)$.
Since $s\log(1/s)\ge c a\log(1/a)$ for $a\le s\le1/2$, we obtain
\begin{equation}\label{cap:set-capacity}
 \mathcal E(h)\ge
 \frac{cA}{(1+\log_+M)^2}\,a\log(1/a).
\end{equation}
This is the required capacity bound, where capacity is the infimum of
$\mathcal E(h)$ over these admissible tests.

For $v\ge0$ supported on a set of measure at most $1/2$, apply
\eqref{cap:set-capacity} to the layers $\phi_j=(v-2^j)_+\wedge2^j$,
$j\in\mathbb Z$, normalized by $2^j$. Writing $a_j=\eta(v\ge2^j)$ and omitting
empty upper level sets gives
\[
 \mathcal E(v)\ge\sum_j\mathcal E(\phi_j)
 \ge\frac{cA}{(1+\log_+M)^2}
       \sum_j4^j a_{j+1}\log(1/a_{j+1}).
\]
For the first inequality, fix an edge and orient its endpoints so that the
increment of $v$ is nonnegative. The increments of all clipped layers are then
nonnegative and sum to that increment. The sum of their squares is at most its
square; summing over edges gives the inequality. Denote the last sum by $S$.
The pointwise dyadic sum $\sum_j4^j\mathbf1_{\{v\ge2^{j+1}\}}$ lies between
$v^2/12$ and $v^2/3$, so $S\ge(\log2)\eta[v^2]/12$. On $2^{j+1}\le v<2^{j+2}$,
the inequalities $\eta[v^2]\ge4^{j+1}a_{j+1}$ and $v^2\le16\cdot4^j$ give
\[
 \log_+\frac{v^2}{\eta[v^2]}
 \le\log\frac4{a_{j+1}}\le3\log\frac1{a_{j+1}}.
\]
Consequently $\operatorname{Ent}_{\eta}(v^2)\le48S$.

Choose a median $b$ of $f$ and apply these estimates to $v_+=(f-b)_+$ and
$v_-=(f-b)_-$. Their supports have measure at most $1/2$, their forms sum to
at most $\mathcal E(f)$, and
\[
 \operatorname{Ent}_{\eta}((f-b)^2)
 \le\operatorname{Ent}_{\eta}(v_+^2)
      +\operatorname{Ent}_{\eta}(v_-^2).
\]
Finally we use an entropy shift inequality of Rothaus type; compare
\cite[Lemma~4]{BR2003}. In the form needed here it is
\begin{equation}\label{cap:entropy-shift}
 \operatorname{Ent}_{\eta}((u+b)^2)
 \le C\operatorname{Ent}_{\eta}(u^2)+C\eta[u^2]
\end{equation}
for every $u$ and every constant $b$. For completeness, normalize
$\eta[u^2]=1$. If $|b|\le2$, use $(u+b)^2\le2u^2+8$ and the positive part of
$s\log s$; its negative part costs at most an absolute constant. If $|b|>2$,
use
\[
 \operatorname{Ent}_{\eta}((u+b)^2)
 \le b^2\eta\!\left[
       \Phi\!\left((1+u/b)^2\right)\right],
 \qquad \Phi(s)=s\log s-s+1.
\]
On $|u|\le|b|/2$ the integrand is at most $Cu^2$ by Taylor's theorem; on its
complement it is at most $Cu^2(1+\log_+u^2)$. This proves
\eqref{cap:entropy-shift}, and hence \eqref{cap:layer-lsi}. The case
$\eta[u^2]=0$ is immediate. \end{proof}

On $\|J\|\le3$, all Gibbs atoms are at least $\exp(-Cn)$.
Lemma~\ref{cap:layer-lemma} therefore turns \eqref{cap:restricted-energy} into
the desired log-Sobolev inequality, after absorbing $(\log n)^2$ by choosing
$r$ smaller than the final exponent allowance. It remains to prove the
restricted estimate.

\subsection{Information and energy under observations}

This is the Gaussian observation form of stochastic localization
\cite{Eldan2013StochasticLocalization,ChenEldan2025Localization}. The
conditional-variance contraction below is also the mechanism of \cite[Lemma~9,
arXiv version~2]{EldanKoehlerZeitouni2022}. We give the identities directly
for the present matrix precisions and keep the stopped-path estimates separate
from any all-field assertion.

Fix a bounded test as in \eqref{cap:restricted-entropy}. For a deterministic
increasing, piecewise smooth matrix precision $B_s$, let $Q_\mu$ and $Q_\nu$
be the observation-path laws obtained from the initial spin laws $\mu$ and
$\nu$, and let $\mathcal F_s$ be the observation filtration. The field has increments
\[
 dh_s=G_sX\,ds+G_s^{1/2}\,dW_s,\qquad G_s=\dot B_s\succeq0.
\]
Let $\mu_s,\nu_s$ be the respective spin posteriors,
$N_s=\mu_s[f^2]$, and write
\[
 l_s=D(\nu_s\Vert\mu_s),\qquad
 H(s)=\mathbb E_{Q_\nu}l_s,\qquad
 m_{\mu,s}=\mu_s[X],\quad m_{\nu,s}=\nu_s[X].
\]
The posterior interaction is $J-B_s$, and
$\nu_s=f^2\mu_s/N_s$. The posterior need not retain the original
half-support or lower entropy-to-log-mass condition. The likelihood and
averaged identities below replace those pointwise restrictions.
These definitions also apply to precision jumps,
by using the corresponding Gaussian channel.

\begin{lemma}[Observation identities]\label{cap:observation-identities}
Along the continuous portions of this path,
\begin{align}
 -\dot H(s)
 &=\tfrac12\mathbb E_{Q_\nu}
       \|m_{\nu,s}-m_{\mu,s}\|_{G_s}^2,
       \label{cap:entropy-identity}\\
 \mathbb E_{Q_\nu}
     \frac{\mathcal D_{\mu_s}(f)}{N_s}
 &\le\frac{\mathcal D_0(f)}N.
       \label{cap:energy-conditioning}
\end{align}
The continuous entropy-loss integral up to a stopping time, together with
the information from preceding precision jumps, gives the relative entropy
of the stopped observation laws. Moreover, for $v\ge0$,
\begin{equation}\label{cap:conditional-entropy-tail}
 Q_\nu(l_s>L+v)\le e^{-v},\qquad
 Q_\nu(E)\le e^L Q_\mu(E)
\end{equation}
for every observation event $E$, including stopped-path events.
\end{lemma}

\begin{proof}
The likelihood ratio of the observation laws through time $s$ is
\[
 \frac{dQ_\nu}{dQ_\mu}\bigg|_{\mathcal F_s}=\frac{N_s}{N}\le e^L.
\]
The
conditional drifts are $G_sm_{\mu,s}$ and $G_sm_{\nu,s}$, with common
innovation covariance $G_s$.  The Gaussian likelihood formula and the
chain rule for relative entropy give
\[
 D(\nu\Vert\mu)
 =D(Q_\nu|_{[0,s]}\Vert Q_\mu|_{[0,s]})+H(s),
\]
and then \eqref{cap:entropy-identity}. The same calculation with the stopped
likelihood proves the stopping assertion; zero-rate directions are omitted.
The finite spin space and bounded finite precision remove any integrability
difficulty.

To prove \eqref{cap:energy-conditioning}, first change the path expectation
from the tilted law to the ordinary law:
\[
 \mathbb E_{Q_\nu}\frac{\mathcal D_{\mu_s}(f)}{N_s}
 =\frac1N\mathbb E_{Q_\mu}\mathcal D_{\mu_s}(f).
\]
The remaining assertion is the conditional-variance
decomposition, applied site by site under the joint ordinary law of the spin
and observations. Conditioning additionally on the observations can only
decrease the expected variance given $X_{-i}$. Finally $f^2\le1$ implies
$N_s/N\le e^L$ and $l_s\le-\log N_s$. Hence
\[
 Q_\nu(l_s>L+v)
 \le\mathbb E_{Q_\mu}
       \left[\frac{N_s}{N}\mathbf1_{\{N_s/N<e^{-v}\}}\right]
 \le e^{-v}.
\]
\end{proof}

We record the ordinary-path estimates used below. Under the original
zero-field Gibbs law observed to precision $B_p$, the quantities
$\|P_pX\|/M_p$ and $\|h_p\|/(p^2M_p)$ have bounded-range tails with
arbitrarily large fixed exponential rates per $k_p$, with the same control for
polynomial moments. The parameter $\alpha$ inferred from the cap field
satisfies
\begin{equation}\label{cap:ordinary-saddle-test}
 Q_\mu\big(|\alpha(h_p)|>\zeta p^2\big)
 \le e^{-a_\zeta k_p}
\end{equation}
for every fixed $\zeta>0$, after reducing the upper cutoff on $p$. These are
quenched estimates on events of probability tending to one. For integration
over logarithmic scale intervals it suffices that the same bounds hold in the
integrated form with a smaller positive exponent.

Here is the additional verification of \eqref{cap:ordinary-saddle-test}. Put
\[
 R_\xi=(2+\xi p^2-J)^{-1},\qquad
 R_{p,\xi}=(2+\xi p^2-J_p)^{-1},
\]
where $\xi>0$ is a sufficiently small fixed constant. Under the
unconditioned Gaussian reference, sample $X^{\rm G}\sim N(0,R_\xi)$ and
$h_p=B_pX^{\rm G}+\eta$, with $\eta\sim N(0,B_p)$ independently. Then
\[
 a_{p,\xi}:=R_{p,\xi}h_p\sim N(0,R_\xi-R_{p,\xi}),\qquad
 \mathbb E\|a_{p,\xi}\|^2=\operatorname{Tr}(R_\xi-R_{p,\xi}).
\]
This expected squared norm differs from
$n-\operatorname{Tr}R_{p,0}$ by $O(\sqrt\xi\,np)+o(np)$.
The covariance has operator norm $O((\xi p^2)^{-1})$ and trace of its square
$O(n/(\sqrt\xi p))$, so a fixed-tolerance deviation has exponent at most
$-c_\zeta\xi k_p$. The one-coordinate lower norm-density estimate for
$X^{\rm G}$ costs $C\xi^{3/2}k_p+O(\log k_p)$ in the logarithm, besides
the density scale $\sqrt{p/n}$. The quadratic tilt retains a spectral band
whose upper norm-density bound has the same scale. Choose $\xi$ small
enough that the $O(\sqrt\xi)$ target error fits the chosen saddle tolerance
and $C\xi^{3/2}<c_\zeta\xi/2$.

To connect this norm test to the saddle, use $h_p\in H_p$ and
$R_{p,0}h_p=(1+\xi)R_{p,\xi}h_p$. The monotone saddle equation and its
scaled inverse bound then place every field with
$|\alpha_p(h_p)|>\zeta p^2$ inside the relaxed quadratic-deviation event.
The preceding conditioning calculation gives that event a positive
spherical rate per $k_p$. The first-moment transfer and zero-field concentration in
Proposition~\ref{cap:static-comparison} prove the claim. The projection and
field tails follow there in the same way, including polynomial moments.
Markov's inequality and Fubini's theorem, summed over geometric scales, give
the stated integrated versions. At a fixed finite number of constant-scale
intervals, time nets and Gaussian maximal inequalities also give whole-path
versions.

The two paths below use at most two fixed endpoint constants $T'>3$,
chosen from their path parameters before $f$. At each endpoint precision
$T'I+J$, the posterior interaction is diagonal and its spin law is product.
On the same kind of disorder event, imposed for this finite collection and
simultaneously for $0\le L\le Cn$,
\begin{equation}\label{cap:product-field-tail}
 Q_\mu\left(\|h_{\rm end}\|_\infty>
                         C\sqrt{\log n+L}\right)
 \le e^{-2L}n^{-10},
\end{equation}
and the omitted tails admit polynomial moment bounds. To check this, condition
on the eigenvalues in the Haar first moment and then on the spin and its
spherical image. The component of $JX$ parallel to $X$ has bounded
coefficient, and its perpendicular component is a uniformly rotated vector of
norm at most $3\sqrt n$. Each coordinate thus has a Gaussian-type tail. Add
the Gaussian observation noise, whose covariance has bounded norm, and take a
coordinate union bound. Increase the rate before the first-moment transfer and
use zero-field partition concentration. A scalar grid in $L$ gives
simultaneity, with stronger tails between grid points. The same calculation
with tail integrals gives the polynomial moment assertion.

\subsection{Retention at small and moderate entropy}\label{cap:entropy-retention}

Fix a small constant $\lambda>0$, to be chosen after the curvature and
ordinary-path tolerances, and first suppose that $L\le\lambda n$. Choose $p_1$
so that, up to fixed endpoint factors,
\begin{equation}\label{cap:first-entropy-scale}
 k_1=np_1^3=\max\{k_{\min},c_0L\},\qquad
 k_{\min}\asymp\sqrt{\log n},
\end{equation}
where $c_0>0$ is a sufficiently small absolute constant.  The coarse
endpoint factors and $c_0$ are chosen before the sharp curvature
tolerances.

The initial observation jump leaves the following amount of entropy:
\begin{equation}\label{cap:jump-retention}
 H(p_1)\ge
 \begin{cases}
  cL,& k_1=c_0L,\\
  e^{-C_0k_1},& k_1=k_{\min}.
 \end{cases}
\end{equation}
To prove the first line, let $Q_0$ be the pure-noise output of the
precision-$B_{p_1}$ channel. Convexity of relative entropy gives $D(Q_\nu\Vert
Q_0)\le\frac12\nu[X^{\mathsf T}B_{p_1}X]$. By Jensen's inequality and
$\mu[X]=0$,
\[
 \log\frac{dQ_\mu}{dQ_0}(h)
 =\log\mu\!\left[e^{h^{\mathsf T}X-X^{\mathsf T}B_{p_1}X/2}\right]
 \ge-\tfrac12\mu[X^{\mathsf T}B_{p_1}X].
\]
Subtract the logarithmic likelihood bound from the pure-noise relative entropy
estimate. The information revealed by the jump is therefore at most half the
sum of these two quadratic moments. The projection tails, transferred also on
an auxiliary dyadic grid, bound this sum by
\[
 Ck_1\left(1+(L/k_1)^{1/3}\right).
\]
For example the $v^6$ tail and $d\nu/d\mu\le e^L$ give
$\nu[\|P_{p_1}X\|^2/M_{p_1}^2]\le C(1+(L/k_1)^{1/3})$. When $k_1=c_0L$, the
displayed loss is at most $C(c_0+c_0^{2/3})L$, which is smaller than $L/4$
after choosing $c_0$; the initial entropy is at least $L/2$.

This argument gives proportional entropy retention when the initial cap
dimension is chosen proportional to L. At the smallest permitted cap that
choice may be unavailable. We then use the support condition instead: the
original law assigns a fixed positive mass to a bounded-projection set that
the tilted law never visits. After a bounded field observation its posterior
mass may be exponentially small in $k_1$, but it remains positive and yields
the weaker entropy bound required in the bottom branch.

For the second line, a set of $\mu$-mass at least $1/4$ lies outside the
support of $f$ and has bounded scaled projection. On a sufficiently large
fixed field ball its posterior mass is at least $e^{-Ck_1}$: the numerator is
bounded below on these projections, and the denominator is bounded above by
integrating the projection tails. The field ball has $Q_\nu$-probability
bounded below because $L\le k_1/c_0$ and its ordinary complement can have
arbitrarily large fixed rate per $k_1$. On this ball $\nu_h$ assigns zero mass
to the indicated complement of the support. Data processing gives
$D(\nu_h\Vert\mu_h)\ge-\log(1-e^{-Ck_1})\ge e^{-C_0k_1}$, after adjusting
$C_0$. These are all coarse constants.

Parameterize the cap observations by $s=\log p$. Their precision rate is
$G=2p^2P_p$. Write $\delta m_h=m_{\nu,h}-m_{\mu,h}$. Suppose that, for
every shift $u\in H_p$ in the neighborhood under consideration,
\[
 P_p\nabla^2\phi_p(h+u)P_p\preceq \frac{C_{\rm ball}}{p^2}P_p.
\]
The required neighborhood contains all such shifts whose radius, in
units $p^2M_p$, is at most
\begin{equation}\label{cap:entropy-shift-radius}
 \left(\frac{2l_h}{k_p C_{\rm ball}}\right)^{1/2}.
\end{equation}
Then the variational formula for relative entropy, tested against linear
functions, gives
\begin{equation}\label{cap:entropy-mean-bound}
 p^2\|P_p\delta m_h\|^2\le2C_{\rm ball}l_h.
\end{equation}
Indeed integrate the Hessian along a shift $u$ to bound its centered log
moment generating function by $C_{\rm ball}\|u\|^2/(2p^2)$, and test in the
direction of $P_p\delta m_h$. The trial shift of length $p\sqrt{2l_h/C_{\rm ball}}$
is available precisely under \eqref{cap:entropy-shift-radius}; the resulting
inequality also bounds the length of the unconstrained optimizer.

We next trade covariance control for a differential inequality on the
remaining entropy. The variational argument requires an entire ball of linear
field shifts, with radius depending on the conditional entropy. A coarse first
interval makes the cap dimension large compared with that entropy. On the
later interval the required shifts fit inside the region of sharp curvature,
apart from errors whose weighted drift contribution is exponentially small.
Keeping those contributions inside the differential inequality is necessary
for the bottom branch, where the retained entropy is itself small.

Here is the bound used for every discarded event. Put
$Y_p(X)=\|P_pX\|/M_p$. Conditional Jensen gives
\[
 p^2\|P_p\delta m_h\|^2
 \le2k_p\bigl(\nu_h[Y_p^2]+\mu_h[Y_p^2]\bigr).
\]
For any observation event $E$, including one depending on $f$, and any
fixed $V>0$, disintegration and the likelihood bound yield
\begin{equation}\label{cap:discarded-drift}
\begin{split}
 \mathbb E_{Q_\nu}\!\left[
 p^2\|P_p\delta m_h\|^2\mathbf1_E\right]
 \le{}&4k_pV^2Q_\nu(E)\\
 &+4k_pe^L\mu\!\left[Y_p^2\mathbf1_{\{Y_p>V\}}\right].
\end{split}
\end{equation}
For example, the tilted posterior term averages to a moment under $\nu$;
the ordinary posterior term uses the path density at most $e^L$.
The projection tail pays the weighted term, so a probability estimate for
$E$ is used together with a moment estimate.

Let $k_2=R_*k_1$ with $R_*$ a sufficiently large fixed constant. On
$[p_1,p_2]$, use the bounded-ball curvature bound of
Theorem~\ref{cap:true-curvature}. Choose $A_\ell\ge A>C_0$, where $C_0$ is
from \eqref{cap:jump-retention}. Discard fields outside a fixed ball and
the event $l_h>L+A_\ell k_p$. Since $L/k_p\le1/c_0$, the latter event has
probability at most $e^{-A_\ell k_p}$ by
\eqref{cap:conditional-entropy-tail}. Enlarge the eventual curvature
constant to satisfy $C_{\rm ball}\ge1$. Before obtaining that constant,
choose the ball to contain shifts of radius
$\sqrt{2(c_0^{-1}+A_\ell)}$ in scaled units from every retained field;
this contains the radius in \eqref{cap:entropy-shift-radius}.
Choose its ordinary field-tail rate and the projection threshold in
\eqref{cap:discarded-drift} larger than $A+1/c_0$.
After paying $e^L$, the discarded drift is at most $Ck_pe^{-Ak_p}$.
The field range and these coarse tail choices precede the sharp tolerances
and $R_*$. Its curvature constant, even if large, causes only a fixed loss
on this interval. Equations \eqref{cap:entropy-identity} and
\eqref{cap:entropy-mean-bound} give
\[
 \frac{dH}{d\log p}\ge-2C_{\rm ball}H-Ck_pe^{-Ak_p}.
\]
The interval has fixed logarithmic length. Integrating this inequality
therefore multiplies the initial lower bound in \eqref{cap:jump-retention} by
a fixed positive factor. The additive error is negligible relative to either
branch of that lower bound.

On $[p_2,p_{\rm top}]$, use instead the curvature $(1+\epsilon)/p^2$ on the
near-zero-parameter region. Choose an ordinary saddle tolerance $\zeta$
strictly inside that region, leaving a field buffer of radius
$r_{\rm buf}$ in units $p^2M_p$. Put
$b_{\rm buf}=(1+\epsilon)r_{\rm buf}^2/2$. The shift radius fits this
buffer whenever $l_h\le b_{\rm buf}k_p$. At these scales
$L/k_p\le(R_*c_0)^{-1}$. Choose $R_*$ so this ratio is below fixed
fractions of both $b_{\rm buf}$ and the rate $a_\zeta$ in
\eqref{cap:ordinary-saddle-test}. Then
\[
 Q_\nu(l_h>b_{\rm buf}k_p)
 \le e^{-(b_{\rm buf}-L/k_p)k_p},\qquad
 Q_\nu(|\alpha_p(h_p)|>\zeta p^2)
 \le e^{-(a_\zeta-L/k_p)k_p}.
\]
Equation~\eqref{cap:discarded-drift}, with another fixed projection
threshold, leaves $Ck_pe^{-ak_p}$ for some fixed $a>0$. Thus
\[
 \frac{dH}{d\log p}\ge-2(1+\epsilon)H-Ck_pe^{-ak_p}.
\]
Increase $R_*$ again so that $aR_*>C_0$ with slack.

For either differential inequality, an exponent $\gamma_H$ and error
$Ck_pe^{-bk_p}$ give
\[
 H(p)\ge\left(\frac{p_a}{p}\right)^{\gamma_H}
 \left[H(p_a)-C\int_{p_a}^{p}
   \left(\frac{u}{p_a}\right)^{\gamma_H}
   k_u e^{-bk_u}\,d\log u\right].
\]
The coarse error starts at $k_1$ with $A>C_0$; the sharp error starts at
$R_*k_1$ with $aR_*>C_0$. Their weighted integrals are therefore negligible
relative to either branch of \eqref{cap:jump-retention}. Solving yields
\begin{equation}\label{cap:retained-entropy}
 H(p_{\rm top})\ge
 c\left(\frac{p_1}{p_{\rm top}}\right)^{2+2\epsilon}
 \begin{cases}
 L,& k_1=c_0L,\\
 e^{-C_0k_1},& k_1=k_{\min}.
 \end{cases}
\end{equation}
Pointwise control of the discarded drift is not needed here. With the
integrated bounds, the integrating factor charges each geometric interval only
a fixed power of $p/p_1$. Its error decreases exponentially in $k_p$, and
$k_1\to\infty$, so their weighted sum is still negligible. Finally choose
$\lambda$ small enough that $R_*c_0\lambda<p_{\rm top}^3$, up to the fixed
endpoint factors. Then $p_2<p_{\rm top}$ for all $L\le\lambda n$.

\subsection{The remaining fixed scales}\label{cap:constant-scales}

We next transfer \eqref{cap:retained-entropy} to a product endpoint, allowing
fixed multiplicative losses. The initialized scale is $p_0$ from
Section~\ref{cap:initialization}, with $t=C_*p_0^2$. Starting from
$p_{\rm top}$, first traverse the finite remaining cap intervals to $p_0$.
Use the actual finer frame for the terminal bridge there. Its increment is
\[
 \Delta_{\rm base}=tI+(1-\beta)J-B_{p_0},
\]
so the total precision is exactly $tI+(1-\beta)J$. The posterior interaction
is $\beta J-tI$, with the same spin law as $\beta J_{\rm off}$.

Now increase scalar precision along
\[
 B_u=uI+(1-\beta)J,\qquad t\le u\le T,
\]
for a sufficiently large fixed $T$. The choice of the single coefficient
$\beta<1$ in Section~\ref{cap:initialization} includes the persistence
requirement for this interval. Namely, fix $T$ and the normalized buffers
for the ordinary critical test after $p_0$ is fixed and before choosing
$\beta$. Adding an independent increment of precision
$(1-\beta)(4I+J)$ to scalar precision $(u-4(1-\beta))I$ gives $B_u$.
The associated common-spin field displacement has normalized size
$O((1-\beta)+\sqrt{1-\beta})$ outside an exponentially rare event.
Time nets and Gaussian maximal bounds preserve the buffered tests along the
interval. Lemma~\ref{cap:endpoint-stability} then supplies a fixed
whole-path rate $e^{-an}$ and a fixed field neighborhood.
The auxiliary removal is performed with this already fixed $\beta$;
it determines $p_{\rm top}$, and $\lambda$ is chosen afterward.

Next use
\[
 B_a=(T+4a)I+(1-\beta+a)J,\qquad 0\le a\le\beta.
\]
Its rate is $4I+J\succ0$, and it ends at
$T'_\beta I+J$ with $T'_\beta=T+4\beta>3$. For $T$ large, uniformly in the underlying spin, the field has large
absolute coordinates except on a small site fraction: the main shift is $TX$,
the other deterministic shifts have bounded normalized norm, and the noise has
squared normalized norm $O(T)$ outside probability $e^{-an}$. At an
approximate root the diagonal $A$ in the stability test is therefore small
outside a small site set. Gaussian nets give a small operator norm for every
sufficiently sparse principal submatrix of a typical GOE matrix. Together with
the full norm bound and the small normalized norm of $A^{1/2}m$, this verifies
the stability test, again on a field neighborhood. While the coefficient of
$J$ belongs to a fixed compact subinterval of $(0,1)$, the classical
covariance bound and the entropy shift argument consequently give
\begin{equation}\label{cap:constant-entropy-ode}
 \dot H\ge-CH-e^{-an},
\end{equation}
after reducing $a$ and $\lambda$ to absorb the path tilt $e^L$.

For sufficiently small interaction coefficients, the covariance is bounded
globally. Indeed add a scalar diagonal to the interaction to obtain a matrix
$Q$ with $0\prec Q\prec I/2$, and introduce a Gaussian latent field with
covariance $Q$. Its posterior log density has Hessian at most
$-Q^{-1}+I\preceq-I$. Brascamp--Lieb bounds the covariance of the
coordinatewise $\tanh$ by $I$, while the conditional spin variances contribute
another $I$. Thus the spin covariance is at most $2I$. The variational
argument now gives $\dot H\ge-CH$ for every field and every $L$, without any
discarded set. Word and residual diagnostics along varying coefficients are
consequently needed only on closed subintervals of $(0,1)$.

The sharp cap estimate stops at $p_{\rm top}$, below the initialization
scales. There are only finitely many intervals left before the classical
field, so bounded amplification and fixed losses suffice there. Starting from
the terminal covariance bound, subtract the two mean-transport identities of
Lemma~\ref{cap:error-transport}. One term transports the difference of
terminal means; the other compares expectations of the bounded terminal
approximation error under two path laws. Pinsker's inequality controls the
latter by the square root of conditional entropy. This propagates the required
quadratic mean estimate backward through the finite list of intervals.

At the terminal classical field, the covariance bound holds on a fixed
normalized neighborhood outside an ordinary event of probability $e^{-an}$.
The linear entropy test fits this neighborhood whenever $l_h\le bn$ for a
fixed $b>0$. Choose $\lambda<\min(a,b)/2$ after these constants. The ordinary
exception costs at most $e^{-(a-\lambda)n}$ under $Q_\nu$, and
\eqref{cap:conditional-entropy-tail} gives
$Q_\nu(l_h>bn)\le e^{-(b-\lambda)n}$. Since both spin means have norm at
most $\sqrt n$, these exceptions have exponentially small quadratic cost.
The neighborhood covariance estimate therefore gives, after reducing the
positive exponent,
\[
 \mathbb E_{Q_\nu}\|m_\nu-m_\mu\|^2\le CH+e^{-an}.
\]
Fix any current time $s$ in one preceding interval and a current field $h$.
Let $Q_{\nu,h}^{\rm cont}$ and $Q_{\mu,h}^{\rm cont}$ be the two laws of the
future observations through the end of that interval, starting from the
posteriors $\nu_h$ and $\mu_h$. They live on the same path space, and data
processing gives
$D(Q_{\nu,h}^{\rm cont}\Vert Q_{\mu,h}^{\rm cont})\le l_h$.
Let $g$ be the gradient of one extended auxiliary potential evolved by heat
convolution, and let $U$ be its linear transport from time $s$ to the interval
end. The same path function $U$ is evaluated under both continuation laws,
and its amplification is bounded by a fixed constant. Subtracting the two
transport identities gives, in every direction $v$ of the current rate subspace,
\begin{align*}
 \langle v,m_\nu-m_\mu\rangle
 ={}&\mathbb E_{Q_{\nu,h}^{\rm cont}}
       \langle Uv,m_{\nu,\rm end}-m_{\mu,\rm end}\rangle\\
 &+(\mathbb E_{Q_{\nu,h}^{\rm cont}}-\mathbb E_{Q_{\mu,h}^{\rm cont}})
       \langle Uv,m_{\mu,\rm end}-g_{\rm end}\rangle.
\end{align*}
On agreement events the second observable is bounded by a constant times
$\|v\|$, by the raw gradient comparison, including at the classical terminal
field. Truncate it there and use the displayed continuation entropy bound and
Pinsker's inequality to bound its expectation difference by
$C\sqrt{l_h}\|v\|$. To bound the complements after averaging, let $W\ge0$
be a terminal cost, such as the squared raw error outside agreement. The
amplification bound lets us dominate the transported exceptional terms by
such costs. Disintegration and the observation likelihood give
\[
 \begin{split}
 \mathbb E_{Q_\nu}\mathbb E_{Q_{\nu,h}^{\rm cont}}W
   &=\mathbb E_{Q_\nu}W
     \le e^L\mathbb E_{Q_\mu}W,\\
 \mathbb E_{Q_\nu}\mathbb E_{Q_{\mu,h}^{\rm cont}}W
   &=\mathbb E_{Q_\mu}\!\left[
      \frac{N_s}{N}\mathbb E_{Q_{\mu,h}^{\rm cont}}W\right]
     \le e^L\mathbb E_{Q_\mu}W.
 \end{split}
\]
The ordinary terminal agreement estimates and their polynomially weighted
versions therefore give exponentially small averaged costs after choosing
$\lambda$ below the finitely many fixed rates. This uses full-path and
current-prefix likelihood bounds, separately from conditional Pinsker.
Apply Cauchy--Schwarz first conditionally on
the current field and then over the tilted observation path. The terminal
entropy satisfies $H_{\rm end}\le H(s)$. If $P_{\rm rate}$ projects onto the
current rate subspace, these bounds give
\[
 \mathbb E_{Q_\nu}\|P_{\rm rate}(m_{\nu,s}-m_{\mu,s})\|^2
 \le C H(s)+e^{-an}.
\]
The transported directions end in the appropriate terminal rate subspace on
cap intervals; for a full-rank bridge $P_{\rm rate}=I$. Since $s$ was
arbitrary, this estimate supplies the differential inequality throughout
the interval.

For the bridge use the actual-refinement TAP extension; for cap intervals use
the auxiliary or true-cap extension supplied by the induction. At each cap
endpoint, \eqref{cap:ordinary-saddle-test} gives agreement outside an
exponentially rare event, and the extension has polynomial gradient moments.
For the bridge, apply Lemma~\ref{cap:marked-terminal-tests}
on the agreement annulus at $p_0$; its ordinary complement is paid by the same
saddle test and polynomial moments. The amplification bound holds from every
intermediate time, while these terminal estimates enter through the averaged
comparison above. No agreement condition at the current field is needed in
the subtraction, since its common approximate gradient cancels.
Iterating through this finite list, even when $p_{\rm
top}$ is interior to one interval, proves \eqref{cap:constant-entropy-ode}.
Its constants may depend on the fixed initialization length. Solving it proves
\eqref{cap:retained-entropy}, up to constants, at the product endpoint.

\subsection{The product estimate}\label{cap:product-estimate}

At precision $T'I+J$ the posterior has independent coordinates with fields
$h_i$. A two-point measure at field $h_i$ has inverse log-Sobolev constant at
most $C(1+|h_i|)$ in conditional-variance normalization. One direct proof
centers the test, uses \eqref{cap:entropy-shift}, and bounds the entropy of
the centered square by its second moment times the logarithm of the inverse
smaller atom; that logarithm is at most $C(1+|h_i|)$. Tensorization and
the posterior entropy identity first give
\[
 l_h\le C(1+\|h\|_\infty)\,
       \frac{\mathcal D_{\mu_h}(f)}{N_h}.
\]
Averaging this bound with \eqref{cap:energy-conditioning} gives
\begin{equation}\label{cap:product-entropy-bound}
 H_{\rm end}\le
 C\sqrt{\log n+L}\,\frac{\mathcal D_0(f)}N+O(n^{-3}).
\end{equation}
To justify the truncation implicit here, apply tensorization on the
event in \eqref{cap:product-field-tail}.  On its complement use
\[
 l_h\le\log\frac1{\min_x\mu_h(x)}
       \le Cn(1+\|h\|_\infty).
\]
The moment version of \eqref{cap:product-field-tail}, followed by the
tilt bound $e^L$, makes its contribution $O(n^{-3})$, uniformly for
$L\le Cn$.

Combining this with \eqref{cap:retained-entropy} proves
\eqref{cap:restricted-energy} for $L\le\lambda n$. Here is the exponent
bookkeeping. In the first case, $p_1^3\asymp L/n$ and $L\gtrsim\sqrt{\log n}$,
so
\[
 \frac{H_{\rm end}}{\sqrt{\log n+L}}
 \ge c n^{-(2+2\epsilon)/3}L
          \frac{L^{(2+2\epsilon)/3}}{\sqrt{\log n+L}}.
\]
The last quotient is bounded below by a negative fixed power of $\log n$ when
$L\le\log n$, and by a constant when $L\ge\log n$. In the second case
$k_1\asymp\sqrt{\log n}$, $p_1^3=k_1/n$, and $L\le C\sqrt{\log n}$; hence
\[
 \frac{H_{\rm end}}{\sqrt{\log n+L}}
 \ge n^{-(2+2\epsilon)/3-o(1)}.
\]
Both lower bounds dominate $n^{-2/3-r}L$ by choosing $\epsilon$ small
relative to $r$.  They also dominate the additive $O(n^{-3})$ error.

\subsection{Large entropy and stopping}\label{cap:large-entropy}

The remaining entropy levels are proportional to the dimension. At these
levels an ordinary rare-event bound cannot simply be multiplied by the initial
density bound. We instead argue under the proposed small energy ratio. The
posterior mean estimate then makes the revealed information small. Relative
entropy of the stopped observation laws transfers the ordinary stopping
probability to the tilted law, and bounded spin means control the loss after
stopping. This retains order-n entropy and conflicts with the product
estimate.

Now suppose $\lambda n\le L\le Cn$.  We prove the stronger estimate
$\mathcal D_0(f)/N>n^{-2/3}L$ by contradiction.  Thus assume
\begin{equation}\label{cap:large-entropy-assumption}
 \frac{\mathcal D_0(f)}N\le n^{-2/3}L\le Cn^{1/3}.
\end{equation}
The objective is to retain a fixed positive multiple of $\lambda n$ until the
product endpoint of this path. Formula \eqref{cap:product-entropy-bound} would
then bound that retained entropy by $O(n^{5/6})+O(n^{-3})$, a contradiction.

We specify the uniformity required at the start. The small-positive-time
statement in Proposition~\ref{cap:tap-expansion}, with the noise interface of
Lemma~\ref{cap:boundary-noise}, uses diagnostics of the centered Gaussian
noise and admits every additional shift of normalized norm at most $Ct$. Its
separate low-band condition is a Gaussian small-ball bound uniform in the
center, not a simultaneous event for all centers. More precisely, on a
disorder event independent of $f$, the noise-only diagnostics fail with
conditional probability $o(1)$, integrated over each fixed interval
$[t_*,t_b]$, and
\begin{equation}\label{cap:small-ball-centers}
 \sup_{z\in\mathbb R^n}
 \mathbb P_g\!\left(
   \|P_t^J(g_t+z)\|<c\sqrt n\,t^{5/4}\right)
 \le e^{-c'nt^{3/2}}.
\end{equation}
Here $P_t^J=\mathbf1_{\{2I-J\le tI\}}=P_{\sqrt t}$ is the cumulative
projection used in Lemma~\ref{cap:boundary-noise}, with $W=J$.
The covariance on a band of rank comparable to $nt^{3/2}$ is comparable to
$tI$; projection onto that band and the Gaussian ball bound prove
\eqref{cap:small-ball-centers}. Translation cannot increase the mass of a
centered Gaussian Euclidean ball. For fixed $J$ the spin is independent of the
observation noise. Since its drift shift has normalized norm $O(t)$ for every
$X\in\{-1,1\}^n$, the deterministic shift allowance and the uniform-in-center
probability give the same $o(1)$ integrated failure bound under every initial
spin law, including $f^2\mu/N$. One first uses Fubini and Markov for the
noise-only diagnostic event to select the disorder event; no union over the
$2^n$ spin centers and no factor $e^L$ occurs in this step. The
exclusion-at-zero and matrix-word events are also selected independently of
$f$.

First fix the large scalar endpoint $T$ and the interaction coefficient below
which covariance is globally bounded. Next choose $t_b>0$ sufficiently small
for Proposition~\ref{cap:tap-expansion}, and obtain the buffered ordinary test
and its failure rate $a>0$ on $[t_b,T]$ and the subsequent weakening interval.
These margins are uniform for all sufficiently small changes of the
interaction coefficient from one. Fix a target integrated drift budget
$\eta_1$ sufficiently small relative to $a\lambda$. Choose $0<t_*\ll t_b$ so
that $Ct_*$ fits that budget, and finally choose $1-\beta'>0$ sufficiently
small relative to these fixed scales and $t_*$.

After fixing $\beta'$, reduce the permitted diagonal enlargement in the
stability test to a positive $\eta_{\rm stab}<(\beta')^{-1}-1$. The
stability implication persists on this smaller diagonal class, and
$(\beta')^2(1+\eta_{\rm stab})^2<1$ supplies the strict margin required by
Proposition~\ref{cap:classical-input}.

Once the resulting stability margins are known, choose the residual tolerance
and finite diagnostic depth in the square-reweighted mean estimate of
Proposition~\ref{cap:classical-input} to realize the target budget. The target
$\eta_1$ is distinct from that proposition's accuracy parameter: the leading
constant multiplying its accuracy may depend on the fixed stability margins.
Its remaining constants may be arbitrarily large but are all fixed before
$n\to\infty$. The factor $n^{-1/2}$ in the posterior estimate will therefore
still make the energy-dependent error vanish under
\eqref{cap:large-entropy-assumption}.

Start directly with precision $t_*I+(1-\beta')J$, increase scalar precision to
$T$, and then weaken the interaction at rate $4I+J$ for duration $\beta'$.
The endpoint is $T'_{\beta'}I+J$ with $T'_{\beta'}=T+4\beta'$, the second
fixed endpoint included in \eqref{cap:product-field-tail}. Positivity of the initial
precision follows from choosing $1-\beta'$ sufficiently small after $t_*$. For comparing this path with
the ordinary critical scalar observation, one may add an independent increment
of precision $(1-\beta')(4I+J)\succ0$ and increase scalar noise time by
$O(1-\beta')$. This produces a small normalized field displacement outside an
exponentially rare event, so the buffered ordinary tests remain valid with
rate $a$ after a fixed decrease.

Let $s_g$ be the onset of the globally bounded-covariance portion. On the
interval $[t_*,t_b]$, the uniform noise and small-ball argument just given
verifies the stability hypotheses with integrated failure tending to zero
under $Q_\nu$. On their validity set the tilted mean estimate from
Proposition~\ref{cap:classical-input}, together with
\eqref{cap:energy-conditioning}, yields an arbitrarily small integrated bound
on $\mathbb E_{Q_\nu}|m_\nu-m_\mu|_n^2$. Indeed the relevant error is bounded
by
\[
 C\eta+\frac{C_\eta}{\sqrt n}
       \left(1+\frac{\mathcal D_{\mu_h}(f)}{N_h}\right),
\]
and its average is $C\eta+O_\eta(n^{-1/6})+O_\eta(n^{-1/2})$ by
\eqref{cap:large-entropy-assumption}. This estimate initially controls a
first moment. Since $|m_\nu-m_\mu|_n\le2$, its square is at most twice its
first power, giving the needed quadratic bound. The integrated failure probability contributes
only $o(1)$ by the same uniform bound.

For times at least $t_b$, use the fixed buffered ordinary stability test,
selected independently of $f$ before observing the path. Let $\tau\le s_g$
be its first violation, taking $\tau=t_b$ if it already fails there
and $\tau=s_g$ if no violation occurs. Up to $\tau$ the same integrated
tilted-mean bound applies deterministically on the diagnostic event. Write
$Q_\nu^\tau,Q_\mu^\tau$ for the laws of the observations stopped at $\tau$.
The initial jump leaks at most $Ct_*n$, by the pure-noise comparison above
without the projection improvement. The stopped version of
\eqref{cap:entropy-identity} then gives
\begin{equation}\label{cap:stopped-path-entropy}
 D(Q_\nu^\tau\Vert Q_\mu^\tau)
 \le (Ct_*+\eta_1+o(1))n,
\end{equation}
where $\eta_1>0$ can be made arbitrarily small by the preceding fixed
accuracy choices.  This includes the possibility of immediate stopping
at $t_b$.

Let $E$ be the event that a violation occurs before the globally
controlled portion.  Its ordinary probability is at most $e^{-an}$.
Data processing for the indicator of $E$ gives
\[
 D(Q_\nu^\tau\Vert Q_\mu^\tau)
 \ge Q_\nu(E)\log\frac1{Q_\mu(E)}-\log2.
\]
Consequently
\begin{equation}\label{cap:stopping-probability}
 Q_\nu(E)\le\frac{Ct_*+\eta_1+o(1)}a.
\end{equation}
On a path that stops, use the spin norm bound for the remaining portion of the
observation. Every mean has norm at most $\sqrt n$, and the precision rate has
bounded norm over a bounded duration. Its entropy-loss integral is
consequently at most $Cn$. Multiplying by the stopping probability bounds the
averaged cost of these paths. By the constant choices,
\eqref{cap:stopped-path-entropy} and \eqref{cap:stopping-probability} make the
total loss before $s_g$ smaller than $\lambda n/4$. Initially $H(0)\ge
L/2\ge\lambda n/2$, so $H(s_g)\ge c\lambda n$. The subsequent global
covariance estimate propagates a fixed fraction to the product endpoint. This
is the claimed contradiction with \eqref{cap:product-entropy-bound}.

All disorder, word, and stability events in this argument have fixed
thresholds selected independently of $f$. The only function-dependent objects
are the tilted law and the entropy level $L$. The proof uses the exponential
ordinary failure rate through stopped-path relative entropy, and therefore
remains valid when $L$ is of order $n$.

\subsection{The log-Sobolev bound}\label{cap:dynamics}

We have proved \eqref{cap:restricted-energy} in both entropy ranges.
Lemma~\ref{cap:layer-lemma} now gives the promised functional inequality.

\begin{theorem}[Critical log-Sobolev upper bound]\label{cap:thm-lsi}
For every fixed $\rho>0$, with probability tending to one over the
Gaussian couplings, the critical zero-field Gibbs measure satisfies
\begin{equation}\label{cap:final-lsi}
 \operatorname{Ent}_\mu(f^2)
 \le n^{2/3+\rho}\mathcal D_0(f)
 \qquad\text{for every }f:\{-1,1\}^n\longrightarrow\mathbb R.
\end{equation}
The form is normalized for the generator
$\mathcal L=\sum_{i=1}^n(Q_i-I)$, where $Q_i$ is conditional expectation
given $X_{-i}$.
\end{theorem}

\begin{proof}
Apply the preceding argument with an exponent allowance strictly
smaller than $\rho$.  On $\|J\|\le3$ the value $M$ in
Lemma~\ref{cap:layer-lemma} is $Cn$.  Its fixed constant and
$(1+\log M)^2$ factor are absorbed in the remaining power of $n$.
\end{proof}

Section~\ref{sec:dynamics} derives the mixing and relaxation bounds from
this inequality and combines them with the companion's lower bounds.
The cap comparison also gives the following direct linear test. It explains
the same worst-start lower exponent with weaker quantifiers.

\begin{remark}[An independent linear lower test]\label{cap:lower-test}
The static comparison also gives a direct, weaker-in-quantifiers route
to the matching worst-start lower exponent.  Fix $0<\epsilon<1/3$ and take
$p=n^{-1/3+\epsilon}$.  In the uncapped zero-field spherical model,
$\|P_pX\|\ge cM_p$ with probability tending to one.  To see this,
compare with the Gaussian of precision $2+\xi p^2-J$.  A spectral band
with gaps comparable to $p^2$ already has this lower bound outside
$e^{-ck_p}$.  The lower norm-density estimate and the bounded density
of the remaining coordinates transfer it through norm conditioning,
when $\xi>0$ is sufficiently small.  Haar first moments and the
zero-field second-moment comparison in
Proposition~\ref{cap:static-comparison} transfer the conclusion to the cube
with high Gibbs probability for typical disorder.  Spin symmetry and
$\operatorname{rank}P_p\le Cnp^3$ imply that some Euclidean unit vector
$a$ has
\[
 \operatorname{Var}_\mu(a^{\mathsf T}X)\ge c p^{-2},
 \qquad
 \mathcal D_0(a^{\mathsf T}X)
 =\sum_i a_i^2\mu[\operatorname{Var}(X_i\mid X_{-i})]\le1.
\]
Thus the rate-one spectral gap is at most $Cp^2$.  Testing a gap
eigenfunction from a state where its absolute value is maximal gives
$t_{\rm mix}(1/4)\ge c p^{-2}$; the discrete gap is smaller by $n$.
Letting the fixed $\epsilon$ decrease proves the same lower exponents.
The realized-start statement and the bounds for every arbitrarily slowly
diverging $M_n$ remain the separate conclusions of
Theorem~\ref{thm:accepted-critical}.
\end{remark}
